\section{An incomplete metric in a big class}
\label{sec:incomplete-divisorial-metric}

We now consider exponential equations $\langle(\theta+\ddc\rho)^n\rangle=\mathrm{e}^{\rho}|\sigma|^{-2}\,\mathrm{d}V$ whose measure has a nonintegrable pole along a divisor, and ask whether the metric $\theta+\ddc\rho$ is complete on the locus where it is smooth and positive. A divisorial part of a solution with minimal singularities can cancel the pole and leave the deleted divisor at finite distance. On the blowup of the projective plane at a point, the resulting metric away from the exceptional divisor is the pullback of the Fubini--Study metric. The class is big but not nef, and the equation has a unique solution of finite energy.

Finite energy refers to the classical closed theory of \cite{SGPT}. For a smooth closed real $(1,1)$-form $\theta$ with big class on a compact K\"ahler surface $Y$, and for $u\in\PSH(Y,\theta)$ with minimal singularities, the Monge--Amp\`ere energy relative to the envelope $V_\theta$ is
\[
 E_\theta(u)=\frac13\sum_{j=0}^2
 \int_Y(u-V_\theta)
 \left\langle(\theta+\ddc u)^j\wedge
                (\theta+\ddc V_\theta)^{2-j}\right\rangle.
\]
For general $u\in\PSH(Y,\theta)$, $E_\theta(u)$ is the decreasing limit of $E_\theta(\max(u,V_\theta-k))$ as $k\to\infty$, and $\mathcal E^1(Y,\theta)$ is the set of potentials with $E_\theta(u)>-\infty$; they have full mass.

\Needspace{18\baselineskip}
\begin{theorem}\label{thm:incomplete-divisorial-metric}\index{Metric completeness!divisorial counterexample}
There exist a smooth projective surface $Y$, a smooth irreducible divisor $E$, a smooth closed real $(1,1)$-form $\theta$ with big class $\alpha$, a Hermitian metric on $\mathcal O_Y(E)$, a defining section $\sigma$ of $\mathcal O_Y(E)$, and a smooth positive volume form $\mathrm{d}V_Y$ with the following properties. The equation
\begin{equation}\label{eq:incomplete-divisorial-equation}
 \left\langle(\theta+\ddc\rho)^2\right\rangle
 =\mathrm{e}^{\rho}|\sigma|^{-2}\,\mathrm{d}V_Y
\end{equation}
has a unique solution $\rho\in\mathcal E^1(Y,\theta)$, and this solution has minimal singularities. Moreover,
\[
 \operatorname{Amp}(\alpha)=Y\setminus E
\]
is precisely the open set on which $\theta+\ddc\rho$ is smooth and strictly positive. The restriction of this metric to every nonempty connected open subset of $Y\setminus E$ is incomplete. In particular, no dense Zariski open subset of $Y\setminus E$ carries a complete induced metric.
\end{theorem}

Let $p=[1:0:0]\in\mathbb P^2$, and let $\omega_2$ be the Fubini--Study form on $\mathbb P^2$, normalized by $\int_{\mathbb P^2}\omega_2^2=1$.

\begin{proof}[Proof of \cref{thm:incomplete-divisorial-metric}]
\emph{The surface and the divisor.}
Use the incidence model
\[
 Y=\{([Z_0:Z_1:Z_2],[W_1:W_2])\in\mathbb P^2\times\mathbb P^1:
           Z_1W_2=Z_2W_1\}.
\]
The projections $\pi\colon Y\to\mathbb P^2$ and $q\colon Y\to\mathbb P^1$ identify $Y$ with the blowup of $\mathbb P^2$ at $p$, with exceptional divisor $E=\pi^{-1}(p)$. Let $\omega_1$ be the Fubini--Study form on $\mathbb P^1$ with the same normalization as $\omega_2$, and write
\[
 \beta=\pi^*\omega_2,\qquad \gamma=q^*\omega_1,
 \qquad \Omega=\beta+\gamma.
\]
The form $\Omega$ is K\"ahler, being the restriction of a product K\"ahler metric.

The linear map $(Z_0,Z_1,Z_2)\mapsto(Z_1,Z_2)$ defines a section $\sigma$ of
\[
 \pi^*\mathcal O_{\mathbb P^2}(1)\otimes q^*\mathcal O_{\mathbb P^1}(-1),
\]
because on $Y$ the vector $(Z_1,Z_2)$ lies in the tautological line $[W_1:W_2]$. The section vanishes simply along $E$, which identifies this line bundle with $\mathcal O_Y(E)$. For the product of the Fubini--Study metric and the metric induced on the tautological bundle, its squared norm is
\[
 s=|\sigma|^2=
 \frac{|Z_1|^2+|Z_2|^2}{|Z_0|^2+|Z_1|^2+|Z_2|^2}.
\]
In a blowup chart $\pi(t,w)=(t,tw)$, $q(t,w)=[1:w]$, one has
\begin{equation}\label{eq:incomplete-divisor-norm}
 s=\frac{|t|^2(1+|w|^2)}{1+|t|^2(1+|w|^2)},
 \qquad \ddc\log s=[E]+\gamma-\beta.
\end{equation}
The current identity follows either from this formula and the second blowup chart or from the Poincar\'e--Lelong formula.

\emph{The big class and its equation.}
Set
\begin{equation}\label{eq:incomplete-current}
 \theta=3\beta-2\gamma,\qquad \rho=2\log s,
 \qquad T=\theta+\ddc\rho=\beta+2[E].
\end{equation}
Thus $\rho\in\PSH(Y,\theta)$. The class $\alpha=\{\theta\}$ is big, because the potential $\psi=(5/2)\log s$ satisfies
\begin{equation}\label{eq:incomplete-strict-current}
 \theta+\ddc\psi
 =\tfrac12(\beta+\gamma)+\tfrac52[E]\geq\tfrac12\Omega.
\end{equation}
It is not nef, since $\int_E\theta=-2$: the restriction of $\beta$ to $E$ vanishes and $q|_E$ is an isomorphism onto $\mathbb P^1$. A nef class has nonnegative degree on every curve, by its defining smooth approximations.

On $Y\setminus E$ define $\mathrm{d}V_Y=\beta^2/s$. To check that it extends across $E$, put
\[
 \mathrm{d}\lambda_{t,w}
 =\tfrac{\mathrm{i}}2\mathrm{d}t\wedge\mathrm{d}\bar t
       \wedge\tfrac{\mathrm{i}}2\mathrm{d}w\wedge\mathrm{d}\bar w,
\]
and let $\mathrm{d}\lambda_{z_1,z_2}$ be the analogous Euclidean volume in affine coordinates of $\mathbb P^2$. The normalized Fubini--Study volume is
\[
 \omega_2^2=\frac{2}{\pi^2}
 \frac{\mathrm{d}\lambda_{z_1,z_2}}{(1+|z_1|^2+|z_2|^2)^3}.
\]
Since the complex Jacobian determinant of $(t,w)\mapsto(t,tw)$ is $t$, we obtain
\begin{align*}
 \beta^2&=\frac{2}{\pi^2}
 \frac{|t|^2\,\mathrm{d}\lambda_{t,w}}
      {[1+|t|^2(1+|w|^2)]^3},\\
 \mathrm{d}V_Y&=\frac{2}{\pi^2}
 \frac{\mathrm{d}\lambda_{t,w}}
      {(1+|w|^2)[1+|t|^2(1+|w|^2)]^2}.
\end{align*}
The second expression is smooth and strictly positive at $t=0$. Interchanging the two homogeneous coordinates gives the extension on the other chart. The two extensions agree off $E$, hence everywhere on their overlap.

The non-pluripolar product of $T$ agrees with $\beta^2$ off $E$ and gives no mass to $E$. The smooth measure $\beta^2$ also gives no mass to $E$, so
\[
 \langle T^2\rangle=\beta^2,
 \qquad
 \mathrm{e}^{\rho}|\sigma|^{-2}\,\mathrm{d}V_Y
   =s^2s^{-1}\frac{\beta^2}{s}=\beta^2.
\]
The right side of \eqref{eq:incomplete-divisorial-equation} is defined on $Y\setminus E$ and extended by zero on $E$, so no product of zero and infinity occurs. This proves the equation and gives $\int_Y\langle T^2\rangle=\int_{\mathbb P^2}\omega_2^2=1$.

\emph{Minimal singularities and energy.}
Let $u\in\PSH(Y,\theta)$. Since
\[
 \beta+\ddc(u-\rho)=\theta+\ddc u\geq0
 \quad\text{on }Y\setminus E,
\]
the function $u-\rho$ on $Y\setminus E$ descends to an $\omega_2$-plurisubharmonic function on $\mathbb P^2\setminus\{p\}$. By \cref{lem:isolated-puncture-extension}, it is bounded above, so $u\leq\rho+C_u$ there. This inequality holds almost everywhere on $Y$, and \cref{lem:qpsh-almost-everywhere-order} makes it hold everywhere. In particular $u=-\infty$ on $E$, and $\rho$ has minimal singularities.

Let $V_\theta$ be the normalized envelope of \cref{sec:obstacle-envelopes}. Since $s\leq1$, we have $\rho\leq0$, and minimality gives
\[
 \rho\leq V_\theta\leq\rho+C.
\]
Thus $\rho-V_\theta$ is bounded off $E$, and the energy is
\[
 E_\theta(\rho)=\frac13\sum_{j=0}^2
 \int_Y(\rho-V_\theta)
 \left\langle(\theta+\ddc\rho)^j\wedge
                (\theta+\ddc V_\theta)^{2-j}\right\rangle.
\]
The mixed measures have finite mass in the closed K\"ahler theory and do not charge $E$. Hence $E_\theta(\rho)>-\infty$, so $\rho\in\mathcal E^1(Y,\theta)$. A potential with minimal singularities has full mass; in this example that mass is $1$.

\emph{Uniqueness.}
Suppose $u\in\mathcal E^1(Y,\theta)$ is another solution, and let $\mathrm{d}\mu=s^{-1}\mathrm{d}V_Y$ on $Y\setminus E$, extended by zero on $E$. This measure is smooth and positive off $E$, but has infinite total mass. The classical comparison principle for full-mass potentials in the closed theory gives
\[
 \int_{\{u<\rho\}}\mathrm{e}^{\rho}\,\mathrm{d}\mu
 \leq\int_{\{u<\rho\}}\mathrm{e}^{u}\,\mathrm{d}\mu.
\]
Both integrals are finite, since the two Monge--Amp\`ere measures have mass $1$. The integrand on the left is strictly larger, which forces $\mu(\{u<\rho\})=0$. Interchanging $u$ and $\rho$ gives $\mu(\{\rho<u\})=0$ as well. Therefore $u=\rho$ almost everywhere on $Y\setminus E$ with respect to a smooth volume form, hence almost everywhere on $Y$, since $E$ has zero volume. By \cref{lem:qpsh-almost-everywhere-order}, $u=\rho$ everywhere.

\emph{The ample locus and incompleteness.}
The K\"ahler current in \eqref{eq:incomplete-strict-current} is smooth off $E$. On the other hand, every $u\in\PSH(Y,\theta)$ satisfies $u\leq\rho+C_u$, so its Lelong numbers at points of $E$ are at least those of $\rho=2\log s$, which are positive. Thus every K\"ahler current in $\alpha$ is singular at every point of $E$, and \cref{def:bc-non-kahler-locus} gives
\[
 \operatorname{Amp}(\alpha)=Y\setminus E_{\mathrm{nK}}(\alpha)
 =Y\setminus E.
\]
The current $T$ equals $\beta$ there, and it is strictly positive because $\pi$ is a biholomorphism away from $E$. Its nonzero divisorial part $2[E]$ prevents smoothness at every point of $E$.

Let $U\subset Y\setminus E$ be a nonempty connected open set, and set $W=\pi(U)$. The map $\pi\colon(U,T|_U)\to(W,\omega_2|_W)$ is an isometry. Choose $x\in W$ and a smooth path $c\colon[0,1]\to\mathbb P^2$ from $x$ to $p$. Its first exit time
\[
 t_0=\min\{t\in[0,1]:c(t)\notin W\}
\]
is positive, and $c([0,t_0))\subset W$. The path has finite Fubini--Study length. For $t_k\nearrow t_0$, the lengths of its tails show that $c(t_k)$ is a Cauchy sequence for the intrinsic metric of $W$. It has no limit in $W$: intrinsic convergence implies ambient convergence, and the ambient limit is $c(t_0)\notin W$. Thus $U$ is incomplete. A dense Zariski open subset of $Y\setminus E$ is the complement of a proper analytic subset in a connected manifold, hence connected, so the last assertion is a special case.
\end{proof}

The cancellation $\mathrm{e}^{\rho}|\sigma|^{-2}=s$ removes the pole, and the divisor $E$, which the metric collapses to the point $p$, stays at finite distance. Thus an exponential equation in a big class can have an incomplete metric even when its finite-energy solution is unique and has minimal singularities. In this example the class is not nef. The next example shows that nefness alone does not restore completeness; there the pole produces a complete cusp, but a different curve lies at finite distance.

\subsection{A nef class with an incomplete null end}
\label{sec:nef-incomplete-cusp}

Nefness does not prevent the failure of completeness. The next construction keeps a complete cusp at the prescribed pole, while a disjoint null curve lies at finite distance.

\begin{theorem}\label{thm:nef-incomplete-cusp}\index{Metric completeness!nef big counterexample}
There are a smooth projective surface $Y$, disjoint smooth divisors $E,D$, a smooth closed semipositive form $\theta$ with nef and big class, a smooth Hermitian norm on $\mathcal O_Y(D)$, its defining section $\sigma$, and a smooth strictly positive volume form $\mathrm{d}V_Y$, such that
\begin{equation}\label{eq:nef-incomplete-cusp-equation}
 \left\langle(\theta+\ddc\rho)^2\right\rangle
 =\mathrm{e}^{\rho}|\sigma|^{-2}\,\mathrm{d}V_Y
\end{equation}
has a unique full-mass solution $\rho\in\mathcal E^1(Y,\theta)$. The solution has zero Lelong numbers and is smooth on $Y\setminus(E\cup D)$, where its current $\Omega=\theta+\ddc\rho$ is strictly positive. The curve $E$ is null for $\{\theta\}$. Every dense Zariski open set on which $\Omega$ is smooth and strictly positive is incomplete, although every path approaching $D$ in $Y\setminus(E\cup D)$ has infinite length.
\end{theorem}

The measure $|\sigma|^{-2}\mathrm{d}V_Y$ has a nonintegrable divisorial pole with smooth strictly positive coefficient. The solution is unbounded at $D$ and does not have ordinary minimal singularities: semipositivity makes the normalized minimal envelope equal to zero. The theorem concerns nefness and bigness; the additional identities required by geometric smoothing remain a separate question.

\begin{lemma}\label{lem:nef-cusp-radial-profile}
There is $F\in C^\infty(\R)$ with $0<p:=F'<1$ and $p'>0$ such that, for $c=1/100$, some $B\in\R$ and $K=(9c^2/8)\mathrm{e}^{-B}>0$,
\begin{align}
 F(t)&=B-3\log(1-c\mathrm{e}^{t/2}),&
 pp'&=K\mathrm{e}^{F+t} &&(t\leq-2),\label{eq:nef-cusp-left-profile}\\
 F(t)&=t+\log H(p(t)),&
 pp'&=\mathrm{e}^{F-t} &&(t\geq2),\label{eq:nef-cusp-right-profile}
\end{align}
where $H(p)=(1-p)^2(1+2p)/6$. At the two ends,
\begin{align*}
 F(t)&=B+3c\mathrm{e}^{t/2}+O(\mathrm{e}^t),&
 p'(t)&\sim(3c/4)\mathrm{e}^{t/2} &&(t\to-\infty),\\
 1-p(t)&\sim2/t,&p'(t)&\sim2/t^2 &&(t\to+\infty),
\end{align*}
and $F(t)-t=-2\log t+\log2+o(1)$ as $t\to+\infty$.
\end{lemma}

\begin{proof}
Put $L(t)=-3\log(1-c\mathrm{e}^{t/2})$ near $(-\infty,-2]$. With $q=c\mathrm{e}^{t/2}$,
\[
 p_L=L'=\frac{3q}{2(1-q)},\qquad
 p_L'=\frac{3q}{4(1-q)^2},\qquad
 p_Lp_L'=\frac{9c^2}{8}\mathrm{e}^{t}(1-c\mathrm{e}^{t/2})^{-3}.
\]
This proves the left identity for every additive constant $B$.

Since $H'(p)=(p-1)p$, the increasing function
\[
 \mathcal T(p)=2+\int_{3/4}^{p}\frac{q}{H(q)}\,\mathrm{d}q
\]
has an inverse $p_R(t)$ for $t\geq2$, extending across $2$, with $p_R(2)=3/4$ and $p_R'=H(p_R)/p_R$. Indeed, its integrand is asymptotic to $2/(1-p)^2$ at $1$. The function $F_R=t+\log H(p_R)$ has derivative $p_R$ and satisfies the right identity. For $q=1-p_R$,
\[
 (1/q)'=\frac{3-2q}{6(1-q)}\longrightarrow\frac12.
\]
Hence $q\sim2/t$, $H(p_R)\sim2/t^2$, and the right-end asymptotics follow.

It remains to interpolate the derivatives. Set $d=p_R(2)-p_L(-2)>0$. Choose disjoint short endpoint collars in $[-2,2]$ and cutoffs $\alpha_L,\alpha_R$ equal to one on smaller collars. For small $\varepsilon>0$,
\[
 Q_0=\alpha_Lp_L'+\alpha_Rp_R'
       +\varepsilon(1-\alpha_L-\alpha_R)
\]
is positive and agrees with the endpoint derivative germs. The collars and $\varepsilon$ can be chosen so that $\int_{-2}^2Q_0<d$. Add a nonnegative smooth interior bump with integral $d-\int_{-2}^2Q_0$. The resulting $Q>0$ gives
\[
 p(t)=p_L(-2)+\int_{-2}^tQ(x)\,\mathrm{d}x.
\]
This matches both endpoint germs, is increasing, and stays in $(0,1)$. Integrating $F'=p$ with $F(2)=F_R(2)$ gives the required $F$ and fixes the left constant $B$. The left expansion follows directly.
\end{proof}

\begin{proof}[Proof of \cref{thm:nef-incomplete-cusp}]
\emph{The background and the pole.}
Use the incidence surface of \cref{thm:incomplete-divisorial-metric}, with blowup projection $\pi:Y\to\mathbb P^2$, second projection $f:Y\to\mathbb P^1$, and exceptional curve $E$ over $[1:0:0]$. Put $D=\pi^{-1}(Z_0=0)$ and $A=\pi^*c_1(\mathcal O(1))$. Then $E\cap D=\varnothing$, $A^2=1$, and $A\cdot E=0$. The class $A$ is nef, since it has a smooth semipositive representative. If $\beta,\gamma$ are the pulled-back Fubini--Study forms, then $\{\gamma\}=A-[E]$ and $\beta+\varepsilon\gamma$ is K\"ahler. Thus
\[
 \frac{\beta+\varepsilon\gamma}{1+\varepsilon}
       +\frac{\varepsilon}{1+\varepsilon}[E]
\]
is a K\"ahler current in $A$, proving bigness.

Choose an even nonnegative smooth function $\zeta$, supported in $[-1,1]$, positive on its interior and of integral one. The smooth convex function
\[
 U(t)=\int_\R\max(0,t-x)\zeta(x)\,\mathrm{d}x
\]
has $U=0$ for $t\leq-1$, $U=t$ for $t\geq1$, and $U',U''\geq0$. On $Z_0\ne0$, write $z=(Z_1/Z_0,Z_2/Z_0)$, $s=|z|^2$, and $t=\log s$. Give the frame $Z_0$ squared norm $\mathrm{e}^{-U(t)}$. It is constant near the affine origin. At infinity, use
\[
 w=Z_0/Z_1,\qquad v=Z_2/Z_1,\qquad h=1+|v|^2,
 \qquad s=h/|w|^2.
\]
The frame $Z_1$ then has squared norm $\mathrm{e}^{-U}/|w|^2=1/h$ near $w=0$. These formulas define a smooth Hermitian metric globally. Pull it back to $\mathcal O_Y(D)$ and let $\sigma=\pi^*Z_0$. Its curvature satisfies
\[
 \theta=\ddc U(t),\quad |\sigma|^2=\mathrm{e}^{-U(t)},\qquad
 \theta=\ddc\log h,\quad |\sigma|^2=|w|^2/h\quad\hbox{near }D.
\]
The radial Hessian proves semipositivity. The form vanishes near $E$ and has rank one near $D$. In particular $V_\theta=0$; the nef class has volume $A^2=1$.

\emph{Extension of the potential.}
On $Y\setminus(E\cup D)\simeq\C^2\setminus\{0\}$, set
\[
 \rho=F(t)-U(t),\qquad \Omega=\ddc F(t).
\]
Its radial Hermitian eigenvalues are $p/s$ and $p'/s$, hence positive. Extend $\rho$ to $E$ by $B$ and to $D$ by $-\infty$.

Near the affine origin the smooth functions $F(\log(|z|^2+\varepsilon))$ are psh and decrease to the continuous extension of $F(\log|z|^2)$. For example, the two Hessian eigenvalues are $p/(s+\varepsilon)$ and $(s p'+\varepsilon p)/(s+\varepsilon)^2$, with $p,p'$ evaluated at $\log(s+\varepsilon)$. Pullback proves positivity near $E$.

Near $D$ put $t_\varepsilon=\log h-\log(|w|^2+\varepsilon)$ and $g(t)=F(t)-t$. The local potentials $g(t_\varepsilon)+\log h$ decrease to $\rho+\log h$. Their Hessians are
\[
 p(t_\varepsilon)\ddc\log h
 +(1-p(t_\varepsilon))\ddc\log(|w|^2+\varepsilon)
 +p'(t_\varepsilon)\frac{\mathrm{i}}{2\pi}
       \partial t_\varepsilon\wedge\bar\partial t_\varepsilon\geq0.
\]
Thus $\rho$ is globally $\theta$-psh. Boundedness at $E$ and the log-log growth at $D$ give zero Lelong numbers everywhere.

\emph{The volume form and full mass.}
With Euclidean area $\mathrm{d}A$, the radial determinant is
\begin{equation}\label{eq:nef-cusp-radial-determinant}
 \Omega^2=\frac{2}{\pi^2}\frac{pp'}{s^2}
                  \mathrm{d}A_{z_1}\mathrm{d}A_{z_2}.
\end{equation}
Define $\mathrm{d}V_Y=\mathrm{e}^{-F}\Omega^2$ off the two divisors. In an exceptional chart $z=(u,uv)$, the real Jacobian is $|u|^2$ and $s=|u|^2h$. The left profile identity gives
\[
 \mathrm{d}V_Y=\frac{2K}{\pi^2h}\,\mathrm{d}A_u\mathrm{d}A_v.
\]
At infinity $z=(1/w,v/w)$ has real Jacobian $|w|^{-6}$. The right identity gives
\[
 \mathrm{d}V_Y=\frac{2}{\pi^2h^3}\,\mathrm{d}A_w\mathrm{d}A_v,
 \qquad
 |\sigma|^{-2}\mathrm{d}V_Y
 =\frac{2}{\pi^2|w|^2h^2}\,\mathrm{d}A_w\mathrm{d}A_v.
\]
The analogous second charts cover each divisor. The extensions are smooth and strictly positive, and agree on overlaps by their equality on the punctured overlaps. The last expression proves the nonintegrable pole with smooth positive coefficient.

On the complement, $\mathrm{e}^{\rho}|\sigma|^{-2}\mathrm{d}V_Y=\Omega^2$. Extend this density by zero on $D$. Both measures give zero mass to the two divisors. Polar integration on $\C^2$ has density $\pi^2s\,\mathrm{d}s$, so
\[
 \int_Y\langle\Omega^2\rangle
 =\int_{-\infty}^{\infty}2pp'\,\mathrm{d}t
 =p(+\infty)^2-p(-\infty)^2=1.
\]
This proves the global equation and full mass.

\emph{Bounded truncations and energy.}
Near $D$ write $\alpha=\ddc t=\ddc\log h$ and $\beta_t=(\mathrm{i}/2\pi)\partial t\wedge\bar\partial t$. They have rank one, and
\[
 \theta=\alpha,\quad \Omega=p\alpha+p'\beta_t,\quad
 \Omega^2=2pp'\alpha\wedge\beta_t,\quad
 \Omega\wedge\theta=p'\alpha\wedge\beta_t,\quad \theta^2=0.
\]
Angular and projective integration of $\alpha\wedge\beta_t$ gives $\mathrm{d}t$. Since $|\rho|=O(\log t)$ and $p'=O(t^{-2})$, its integrals against $\Omega^2$ and $\Omega\wedge\theta$ are finite there. Near $E$ the potential is bounded and $\theta=0$; the remaining region is compact with smooth data.

Choose $M$ so that $\phi=\rho-M\leq0$, and set $\phi_k=\max(\phi,-k)$. For large $k$ the truncation changes only the cusp. There $\phi=\log H(p(t))-M$ is strictly decreasing. If $\phi(T_k)=-k$, then $k=2\log T_k+O(1)$. The punctured radial primitive is
\[
 G_k(t)=\max(F(t)-M,t-k).
\]
Its genuine local potential is $G_k(t)+\log|w|^2$, equal to $\log h-k$ near $D$. This local frame term prevents a spurious divisor mass. The derivative jumps upward from $p(T_k)$ to $1$.

Convolving $G_k$ in $t$ with an even nonnegative smooth kernel gives convex primitives decreasing to $G_k$ as the scale decreases. The local bounded potentials converge uniformly on compact sets, so their Bedford--Taylor products converge. Integrating the smooth radial square and mixed measures through the transition gives respectively the shell masses
\[
 1-p(T_k)^2,\qquad 1-p(T_k).
\]
Beyond the transition the primitives are linear and both top-degree measures vanish. On the shell $\phi_k=-k$, so the absolute energy contributions are $k(1-p(T_k)^2)$ and $k(1-p(T_k))$. Both tend to zero, since $1-p(T_k)\sim2/T_k$. The interior contributions converge absolutely by the preceding bounds. At $E$, bounded potentials have no mass on the curve, consistent with the zero radial boundary mass $p(-\infty)^2=0$. Consequently
\[
 E_\theta(\phi_k)=\frac13\sum_{j=0}^2
       \int_Y\phi_k(\theta+\ddc\phi_k)^j\wedge\theta^{2-j}
\]
has a finite limit and is uniformly bounded below. This proves $\rho\in\mathcal E^1(Y,\theta)$.

Uniqueness follows by the full-mass comparison argument in \cref{thm:incomplete-divisorial-metric}, with $\mathrm{d}\mu=|\sigma|^{-2}\mathrm{d}V_Y$ off $D$, extended by zero on $D$. Its total mass is infinite, but each exponential density is a finite measure of mass $1$. Comparison on the two strict sublevel sets forces equality almost everywhere, and \cref{lem:qpsh-almost-everywhere-order} gives equality of the potentials everywhere.

\emph{The two metric ends.}
Fix $v=v_0$ in an exceptional chart and write $h_0=1+|v_0|^2$. The normal restriction is
\[
 \Omega|_{v=v_0}=\frac{\mathrm{i}}{2\pi}
       \frac{p'(\log(|u|^2h_0))}{|u|^2}
       \,\mathrm{d}u\wedge\mathrm{d}\bar u.
\]
The coefficient is asymptotic to $(3c\sqrt{h_0}/4)|u|^{-1}$, so the current has no smooth extension at any point of $E$. The radial path $u=r>0$ has finite remaining length, bounded by $C\int_0^{r_0}r^{-1/2}\,\mathrm{d}r$. At $D$ the normal coefficient is asymptotic to $2/(|w|^2t^2)$ and also diverges. Thus every smooth positive regular locus excludes both curves.

Let $W$ be any dense Zariski open positive regular locus. The proper analytic set $Y\setminus W$ has finitely many curve and point components, by local finiteness of irreducible components and compactness. Distinct curves meet $E$ in finitely many points. Choose a point of $E$ avoiding all these intersections and point components. In a sufficiently small chart the complement is exactly $E$, so a tail of the preceding radial path lies in $W$. Its points approaching $E$ form an intrinsic Cauchy sequence with no limit in $W$: the smooth positive metric induces its manifold topology, while the only ambient limit is the removed point. Hence $W$ is incomplete.

Finally, on the cusp $\Omega=p\ddc t+p'\,\mathrm{d}t\wedge\dc t$, with both terms nonnegative and $p'\geq c_0/t^2$ for large $t$. The length of any path with $t\to\infty$ is bounded below by a positive multiple of $\int|\mathrm{d}t|/t$, which diverges. The missing finite-distance end is therefore the null curve disjoint from the pole.
\end{proof}


\subsection{An incomplete metric with compatible flow data}
\label{sec:flow-compatible-incomplete-metric}

The auxiliary divisorial equation can be incomplete even when its class, divisor and volume come from the prescribed initial data of a geometric smoothing problem.

Let $(X,\omega)$ be compact K\"ahler, with an $\omega$-psh initial potential $\phi_0$ having analytic singularities. The compatible data on a log resolution $\pi:Y\to X$ satisfy
\begin{gather}
 \pi^*\phi_0=\sum_jm_j\log|\sigma_j|^2+\ell_0,\qquad
 \pi^*(\omega+\ddc\phi_0)=\sum_jm_j[D_j]+R_0,
 \label{eq:compatible-initial-resolution}\\
 K_Y=\pi^*K_X+\sum_jb_jD_j,\qquad
 \pi^*\mathrm{d}V_X=\prod_j|\sigma_j|^{2b_j}\mathrm{d}V_Y.
 \label{eq:compatible-volume-resolution}
\end{gather}
Here $D=\sum_jD_j$ is SNC, the defining sections have smooth Hermitian squared norms, $\ell_0$ is smooth, and $R_0$ is smooth closed semipositive. If $\Theta_j$ denotes the curvature of the section norm, put
\[
 \lambda=\max_j\frac{m_j}{1+b_j},\qquad
 \Theta=\sum_j(1+b_j)\Theta_j,\qquad
 \theta=\lambda\Theta+R_0,\qquad
 |\sigma|^2=\prod_j|\sigma_j|^2.
\]
The number $\lambda$ is the inverse exponential-integrability index of $\phi_0$. The metric in the next theorem is the solution of the auxiliary elliptic equation at this background.

\begin{theorem}\label{thm:flow-compatible-incomplete-metric}\index{Metric completeness!compatible initial data}
There are compatible data \eqref{eq:compatible-initial-resolution}--\eqref{eq:compatible-volume-resolution} on smooth projective surfaces, with $\mathrm{d}V_X=\omega^2$, $R_0$ K\"ahler and $D$ a disjoint union of two smooth curves, such that
\begin{equation}\label{eq:compatible-incomplete-equation}
 \langle(\theta+\ddc\rho)^2\rangle
 =\mathrm{e}^{\rho}|\sigma|^{-2}\,\mathrm{d}V_Y
\end{equation}
has a unique full-mass solution in $\mathcal E^1(Y,\theta)$. It has the Lelong numbers of a minimal potential, but every dense Zariski open positive regular locus of its current is incomplete. The unweighted pole is nonintegrable, with smooth strictly positive coefficient across both components of $D$.
\end{theorem}

The auxiliary class in this example is non-nef. Thus the theorem and \cref{thm:nef-incomplete-cusp} answer the two separate completeness assertions negatively. Completeness under simultaneous nefness and these compatibility identities remains open. Neither example identifies its auxiliary metric with the evolving maximal flow.

\Needspace{12\baselineskip}
\begin{lemma}\label{lem:fixed-volume-curve-cusp}
On $\mathbb P^1$, let $z=Z_1/Z_0$, $r=|z|^2$, $\alpha=\ddc\log(1+r)$ and $h=r/(1+r)$, with $\int\alpha=1$. For every $c>2$ there is a rotational $c\alpha$-psh function $u$, smooth off $0$, such that
\begin{equation}\label{eq:fixed-volume-curve-cusp}
 g_1:=c\alpha+\ddc u=\mathrm{e}^{u}h^{-1}c\alpha.
\end{equation}
It has zero Lelong number at $0$, mass $\int g_1=c$, and
\[
 \int_{\mathbb P^1}|u|\alpha+\int_{\mathbb P^1}|u|g_1<\infty.
\]
For $s=-\log r\to\infty$, $t=\log r$ and some $A\in\R$,
\[
 \begin{aligned}
 u&=\log(2/c)-2\log(s+A)+O(\mathrm{e}^{-s}),\\
 u_t&=\frac{2}{s+A}+O(\mathrm{e}^{-s}),\qquad
 u_{tt}=\frac{2}{(s+A)^2}+O(\mathrm{e}^{-s}).
 \end{aligned}
\]
\end{lemma}

\begin{proof}
\emph{Barriers and the radial inverse.}
Set $L=1-\log h=1+\log(1+1/r)$ and $w=-2\log L$. Differentiation gives
\[
 \ddc w=\left(-\frac2L+\frac2{rL^2}\right)\alpha,\qquad
 Q(r):=\frac{c\alpha+\ddc w}{\mathrm{e}^{w}h^{-1}c\alpha}
 =\frac{crL^2-2rL+2}{c(1+r)}.
\]
The function $Q$ is positive, since $L\geq1$ and $c>2$. Its limits are $2/c$ at $0$ and $1-2/c$ at infinity, so $0<q_-\leq Q\leq q_+<\infty$. Choose $a_0\leq\min(0,\log q_-)$ and $b_0\geq\max(0,\log q_+)$. Then $l=w+a_0$ and $H=w+b_0$ satisfy
\[
 c\alpha+\ddc l\geq\mathrm{e}^{l}h^{-1}c\alpha,\qquad
 c\alpha+\ddc H\leq\mathrm{e}^{H}h^{-1}c\alpha,\qquad l\leq w\leq H.
\]

Exhaust $\mathbb P^1\setminus\{0\}$ by disks $|\zeta|<R$ in $\zeta=1/z$, with boundary value $w$. For $q=|\zeta|^2$, the equation becomes
\[
 \Delta u=f(q)\mathrm{e}^{u}-g(q),\qquad
 f(q)=\frac{4c}{1+q},\qquad g(q)=\frac{4c}{(1+q)^2},
\]
where $\Delta$ is the Euclidean real Laplacian.

For $k>0$, the regular radial solution $v=V(q)$ of $(-\Delta+k)v=F(q)$ satisfies
\begin{equation}\label{eq:radial-cusp-volterra}
 V(q)=A_0+\frac14\int_0^q\log(q/t)\bigl(kV(t)-F(t)\bigr)\,\mathrm{d}t.
\end{equation}
The $j$-fold iterate of the integral kernel on $1$ is $q^j/(4^j(j!)^2)$, so successive integration converges uniformly on every finite interval. The coefficient of $A_0$ is $\sum_{j\geq0}(kq/4)^j/(j!)^2>0$. Hence a unique $A_0$ gives any prescribed value at $q=R^2$. Moreover,
\[
 V'(q)=\frac14\int_0^1\bigl(kV(qt)-F(qt)\bigr)\,\mathrm{d}t.
\]
This formula gives smoothness at $q=0$ by induction. The maximum principle for $-\Delta+k$ proves uniqueness and order preservation of the radial inverse.

Choose $k\geq\sup f\mathrm{e}^{H}$ on the fixed disk. Invert $-\Delta+k$ with boundary value $w$ and right side $kv-f\mathrm{e}^{v}+g$, which is nondecreasing on $[l,H]$. The barrier inequalities give $l\leq T(l)$ and $T(H)\leq H$. Iteration from $l$ therefore increases and stays below $H$. The displayed derivative formula gives a uniform first-derivative bound on the fixed interval. Compactness and monotonicity give uniform convergence; passing to \eqref{eq:radial-cusp-volterra} proves the equation, and the averaged derivative formula bootstraps the limit to smoothness.

The same barriers bound the solutions on every fixed reciprocal disk as $R\to\infty$. For an exact nonlinear solution $U(q)$,
\[
 U'(q)=\frac14\int_0^1\bigl(f(qt)\mathrm{e}^{U(qt)}-g(qt)\bigr)\,\mathrm{d}t.
\]
This gives all derivative bounds on compact intervals by induction. A diagonal limit is smooth on the reciprocal plane, including $z=\infty$, and has $u=-2\log(-\log r)+O(1)$ at $0$.

\emph{The asymptotics, mass and energy.}
The radial equation is
\[
 u_{tt}=\frac{c\mathrm{e}^{u}}{1+\mathrm{e}^t}
              -\frac{c\mathrm{e}^t}{(1+\mathrm{e}^t)^2}.
\]
Put $v(t)=u(t)+c\log(1+\mathrm{e}^t)$ and $V(s)=v(-s)$. Then
\[
 V''=k_0(s)\mathrm{e}^{V},\qquad
 k_0(s)=c(1+\mathrm{e}^{-s})^{-(c+1)},\qquad
 V=-2\log s+O(1).
\]
Convexity and this sublinear decay imply $V'<0$ and $V'\to0$ at infinity. The function $J=(V')^2-2k_0\mathrm{e}^{V}$ satisfies $J'=-2k_0'\mathrm{e}^{V}$ and tends to zero. Since $k_0'=O(\mathrm{e}^{-s})$ and $\mathrm{e}^{V}\asymp s^{-2}$,
\[
 J(s)=2\int_s^\infty k_0'(\tau)\mathrm{e}^{V(\tau)}\,\mathrm{d}\tau
       =O(\mathrm{e}^{-s}s^{-2}).
\]
Thus $(\mathrm{e}^{-V/2})'=\sqrt{c/2}(1+O(\mathrm{e}^{-s}))$, and integration gives $\mathrm{e}^{-V/2}=\sqrt{c/2}(s+A)+O(\mathrm{e}^{-s})$. This identity and the differential equation give all three asserted asymptotics.

The local potential $u+c\log(1+r)$ extends subharmonically by \cref{lem:qpsh-removable-analytic}. Its radial slope tends to zero at the puncture, so the extension adds no point mass. The density is asymptotic to
\[
 \frac{\mathrm{i}}{2\pi}\frac{2}{r s^2}
          \,\mathrm{d}z\wedge\mathrm{d}\bar z.
\]
Distributional Stokes gives mass $c$. The energy integral at the cusp is bounded by $C\int^\infty(\log s)/s^2\,\mathrm{d}s$; the integral against the smooth form $\alpha$ is finite as well.
\end{proof}

\Needspace{12\baselineskip}
\begin{proof}[\cref{thm:flow-compatible-incomplete-metric}]
\emph{The initial data and the resolution.}
Take $X=\mathbb P^1_Z\times\mathbb P^1_W$, with normalized Fubini--Study forms $\alpha,\beta$, and put
\[
 \omega=6(\alpha+\beta),\qquad \mathrm{d}V_X=\omega^2=72\alpha\wedge\beta.
\]
Let $H=(Z_1=0)$ and $p=([0:1],[1:0])=(\infty,0)$. Define
\[
 \begin{gathered}
 h_H=\frac{|Z_1|^2}{|Z_0|^2+|Z_1|^2},\qquad
 \eta=\frac{|Z_0|^2}{|Z_0|^2+|Z_1|^2}
              +\frac{|W_1|^2}{|W_0|^2+|W_1|^2},\\
 \phi_0=2\log h_H+2\log\eta.
 \end{gathered}
\]
The zero sets are $H$ and $\{p\}$, which are disjoint. Blow up $p$, write $E$ for the exceptional curve, and set $H_Y=\pi^{-1}H$, $D=H_Y+E$. The projective blowup is smooth, compact, connected and K\"ahler, and is an isomorphism off the initial singular locus.

Use $|\sigma_H|^2=\pi^*h_H$ and $|\sigma_E|^2=\pi^*\eta$. Near $p$, put $x=Z_0/Z_1$, $y=W_1/W_0$. In the chart $x=a$, $y=ab$,
\[
 \pi^*\eta=|a|^2A(a,b),\qquad
 A(a,b)=\frac1{1+|a|^2}+\frac{|b|^2}{1+|ab|^2}>0.
\]
This defines a smooth Hermitian metric on $\mathcal O_Y(E)$ across $E$. The second chart is obtained by interchanging $x,y$, and the metric factors satisfy the transition rule.

The sections $Z_0W_0$, $\sqrt2Z_0W_1$, $Z_1W_1$ of $\mathcal O(1,1)$ have summed squared norm $\eta$ and base ideal $(x,\sqrt2xy,y)=(x,y)$. Dividing by $\sigma_E$ gives a morphism $F:Y\to\mathbb P^2$ with
\begin{equation}\label{eq:compatible-curvature-map}
 F^*\gamma=\pi^*(\alpha+\beta)-\Theta_E,\qquad
 \Theta_H=\pi^*\alpha,
\end{equation}
where $\gamma$ is normalized Fubini--Study. In the first chart $F=[1:\sqrt2ab:b]$, so $F|_E$ is an isomorphism onto a line.

Near $p$,
\[
 \eta=\frac{|x|^2+|y|^2+2|xy|^2}{(1+|x|^2)(1+|y|^2)}.
\]
Thus $\phi_0$ has analytic singularities; in the convention $\frac1{2m}\log\sum|f_j|^2$, the coefficient $2$ is obtained by taking the fourth power of the sum of squares and expanding it as a sum of squared holomorphic monomials. The exact resolved identities are
\begin{align}\label{eq:compatible-kahler-residual}
 \pi^*\phi_0&=2\log|\sigma_H|^2+2\log|\sigma_E|^2,\\
 R_0&=\pi^*\omega-2\Theta_H-2\Theta_E
       =2\pi^*\alpha+4\pi^*\beta+2F^*\gamma.\notag
\end{align}
The first two terms of $R_0$ are positive away from $E$. On $E$, their kernel is $TE$, and $F^*\gamma$ is positive on $TE$. Hence $R_0$ is K\"ahler. The resolved current is $2[H_Y]+2[E]+R_0$, which proves $\omega$-psh positivity away from $p$; the bounded-above extension across $p$ gives it everywhere. We have $m_H=m_E=2$ and $\ell_0=0$.

The blowup Jacobian is $a$, so $b_E=1$, $b_H=0$. Direct exponential integrability requires $q<1/2$ at $H$ and $q<1$ at $p$. Consequently $\lambda=2=\max\{2,1\}$, and
\begin{equation}\label{eq:compatible-auxiliary-background}
 \Theta=\Theta_H+2\Theta_E,\qquad
 \theta=2\Theta+R_0=\pi^*\omega+2\Theta_E.
\end{equation}
Since $\Theta_E|_E=-\omega_{\mathrm{FS}}$, the class has intersection $-2$ with $E$ and is non-nef. It is big: the current associated to $3\log|\sigma_E|^2$ is
\[
 (\pi^*\omega-\Theta_E)+3[E],\qquad
 \pi^*\omega-\Theta_E=5\pi^*(\alpha+\beta)+F^*\gamma>0.
\]

\emph{The exact prescribed volume.}
Define $\mathrm{d}V_Y=\pi^*\omega^2/|\sigma_E|^2$ off $E$. In the chart above,
\[
 \mathrm{d}V_Y=\frac{72\,\mathrm{d}A_a\mathrm{d}A_b}
 {\pi^2A(a,b)(1+|a|^2)^2(1+|ab|^2)^2}.
\]
The cancellation of the Jacobian $|a|^2$ is exact. The second chart gives the extension on the rest of $E$, and the extensions agree on overlaps. Thus the volume is globally smooth and strictly positive, with $\pi^*\mathrm{d}V_X=|\sigma_E|^2\mathrm{d}V_Y$ as required. Since $H_Y$ and $E$ are disjoint, $|\sigma_H\sigma_E|^{-2}\mathrm{d}V_Y$ has a nonintegrable ordinary pole with smooth positive coefficient at each component.

\emph{The solution and the minimal reference.}
Take $u$ from \cref{lem:fixed-volume-curve-cusp} with $c=6$, pulled to the first factor of $X$, and put $g_X=(6\alpha+\ddc u)+6\beta$. It is smooth and strictly positive near $p$, and
\[
 g_X^2=\mathrm{e}^{u}h_H^{-1}\omega^2.
\]
Set $P=2\log|\sigma_E|^2$ and $\rho=P+\pi^*u$. Then
\[
 \theta+\ddc P=2[E]+\pi^*\omega,\qquad
 \theta+\ddc\rho=2[E]+\pi^*g_X.
\]
Off $D$ the equation follows from the exact identity
\[
 \mathrm{e}^{\rho}|\sigma_H\sigma_E|^{-2}\mathrm{d}V_Y
 =\mathrm{e}^{\pi^*u}|\sigma_H|^{-2}|\sigma_E|^2\mathrm{d}V_Y
 =\pi^*g_X^2.
\]
Both measures extend by zero on $D$, proving the global non-pluripolar equation.

For any $\theta$-psh potential $\psi$, the function $\psi-P$ descends off $E$ to an $\omega$-psh function on $X\setminus\{p\}$. By \cref{lem:isolated-puncture-extension}, it extends and is bounded above on $X$. Thus $\psi\leq P+C_\psi$ off $E$, and \cref{lem:qpsh-almost-everywhere-order} gives the inequality everywhere. Since $P$ itself is $\theta$-psh, it has minimal singularities.

Its current has non-pluripolar mass $\int_X\omega^2=72$, the volume of the class. In contrast, the ordinary cup-product square is $72+4E^2=68$. The cusp has mass $6$ and no atom, so the current at $\rho$ also has mass $\int_Xg_X^2=72$. Thus $\rho$ has full mass. Near $E$, $\pi^*u$ is smooth, giving the Lelong numbers of $P$; on $H_Y$ the log-log cusp has zero Lelong number. The actual solution is not of ordinary minimal type, because $u$ is unbounded below at $H$.

\emph{Relative energy and uniqueness.}
Choose $M$ with $u-M\leq0$ and use the minimal reference $P+M$. Its canonical approximants are
\[
 \rho_k=P+M+\pi^*u_k,\qquad u_k=\max(u-M,-k).
\]
Their relative difference off $E$ is $\pi^*u_k$. On the curve let $b(t)=6\log(1+\mathrm{e}^t)$ and $V(t)=b(t)+u(t)$. For large $k$ there is a unique cusp level $t=T_k$ with $u(T_k)-M=-k$. If $S_k=-T_k$, the asymptotics give $S_k\asymp\mathrm{e}^{k/2}$.

The truncated curve primitive is $b+u_k=\max(V-M,b-k)$. Below $T_k$ its mass is $b'(T_k)$; at $T_k$ the derivative jumps upward by $u_t(T_k)$. Smooth convex approximation and bounded Bedford--Taylor continuity therefore give total cutoff mass $V'(T_k)=O(S_k^{-1})$. On that region $u_k=-k$, so its absolute energy contribution is
\[
 kV'(T_k)=O(k\mathrm{e}^{-k/2})\longrightarrow0.
\]
Outside it the original density and potential are unchanged, and $\int|u|g_1$ is finite. Consequently $\int|u_k|g_{1,k}$ has a finite limit, where $g_{1,k}=6\alpha+\ddc u_k$. Also $\int|u_k|\alpha$ converges by domination.

The mixed measures for the product current and its reference are
\[
 \omega^2=72\alpha\wedge\beta,\quad
 (g_{1,k}+6\beta)\wedge\omega
     =6g_{1,k}\wedge\beta+36\alpha\wedge\beta,\quad
 (g_{1,k}+6\beta)^2=12g_{1,k}\wedge\beta.
\]
Their pullbacks are precisely the mixed non-pluripolar measures upstairs, ignoring $E$. All three canonical relative energy integrals therefore have finite limits. This proves $\rho\in\mathcal E^1(Y,\theta)$. Replacing $P+M$ by the normalized minimal reference changes it by a bounded amount and preserves this finite-energy class.

Uniqueness follows by the same full-mass comparison argument as in \cref{thm:incomplete-divisorial-metric}. The unweighted pole measure has infinite mass, but the two exponential densities being compared have mass $72$. Comparison on the strict sublevel sets forces equality almost everywhere, and then equality of the canonical quasi-psh representatives.

\emph{The missing metric boundary.}
The positive smooth locus of the current is $Y\setminus(H_Y\cup E)$. The cusp coefficient diverges at $H_Y$, and the nonzero current $2[E]$ excludes smoothness at $E$. The pullback metric part is also degenerate along $TE$.

For any dense Zariski open positive regular locus $U$, choose a point of $E$ avoiding every other component of its analytic complement, as in \cref{thm:nef-incomplete-cusp}. In a small chart the complement is exactly $E$. A transverse path $a=\xi>0$, $b=b_0$ lies in $U$ for small $\xi$ and approaches that point. Since $g_X$ is smooth near $p$ and $\pi$ is smooth, the coefficients of $\pi^*g_X$ are bounded in the chart. The path has bounded speed and finite remaining length. Its tails give an intrinsic Cauchy sequence with no limit in $U$, since its only ambient limit is the removed point and a smooth positive Riemannian metric induces the manifold topology. Hence every such $U$ is incomplete.
\end{proof}

